% !TeX root = ../../Book3.tex
\chapter*{后记}
\addcontentsline{toc}{chapter}{后记}
\markboth{后记}{后记}

本卷以两个方程开始：\(x=0\) 与 \(x^2=0\)。它们在同一域 \(k\) 中有同一个零点，相应商环却不同。在 \(k[x]/(x^2)\) 中，\(x\) 的类仍然存在，只是平方为零。读到这里，这个差别已经不再是孤立的例外。幂零元、嵌入分量、交点的长度和纤维的厚度，都要求我们保留点集没有记录的信息。理想与商环使这些信息能够参与运算，素谱及其结构层则使它们能够在空间上被辨认。

\ProseParagraphBreak
从方程转到环，并没有让具体计算失去意义。消去仍然要回答哪些坐标可以舍去，理想分解仍然要辨认不同分量，有限商的基仍然要逐项算出。变化在于，计算的输出需要说明自己代表什么。零点集相同不保证理想相同，投影的闭包不一定就是投影，软件列出的多项式也须由生成元表示与首项条件核验。一个可使用的答案，应当连同它所在的环、所保留的关系和可以检查的等式一起给出。

\ProseParagraphBreak
局部化、完备化与伴随分次环使我们能够提出更细的问题。局部化允许特定元素成为分母；完备化加入所有阶相容的形式极限；伴随分次环记录逐层留下的最低阶信息。它们各有用途，也各自改变了观察的对象。同样一组方程写在分次环、局部环和完备局部环中，不能因此被当作同一个环；茎与纤维之间也不能省去扩张标量这一步。有限截断可以用于计算，但必须另有理由，才能让这个有限答案决定完整长度或完整关系。

\ProseParagraphBreak
维数以后还需要深度与同调。正则序列检验逐次取商是否遇到零因子，自由消解记录生成元之间的关系以及关系之间的关系，\(\operatorname{Tor}\) 与 \(\operatorname{Ext}\) 使这些问题有了可比较的次数。Cohen--Macaulay 性要求深度达到维数，正规性则结合余维一的正则性与 Serre 深度条件；它们不能排成一条性质逐渐变弱的序列。具体的消解与反例，是理解这些区别的依据。对非 Cohen--Macaulay 环，典范模还不足以概括全部对偶信息，必须保留多个次数及联系它们的复形。

\ProseParagraphBreak
从局部返回整体，始终需要条件。一个理想处处局部主生成，仍可能没有全局生成元；一个模局部自由，仍可能需要非平凡的过渡同构。层把这些过渡资料纳入对象，概形把各点附近的函数环组织起来，上同调则记录某些局部选择不能成为全局选择的障碍。于是，局部方法的作用既包括建立整体结论，也包括说明结论不能成立的原因。理想类、黏合资料和同调群，使这些障碍能够被准确陈述，而不只是归结为计算失败。

\ProseParagraphBreak
本卷已经给出层、仿射模层对应、Proj、纤维积和若干射影上同调的具体构造。继续研究一般相对几何、导出范畴及更广的对偶理论，还需要新的准备。附录中的数域、奇点、微分消去、局部计算与实正性，也分别指向各自的专题；其中所采用的系数对象存在性、算法分解和解析表示等深层结果，保留了引用与证明范围。能够核验一个模型，并不等于已经获得整个领域的理论，但模型应当使下一步的问题具体可见。

\ProseParagraphBreak
我希望读者留下的习惯，是在使用一个结论之前，先确认眼前的对象和映射：底环是什么，允许哪些分母，是否经过完备化或基变换，有限性用在何处，所得证书究竟证明了什么。再回到 \(x=0\) 与 \(x^2=0\)，我们已经能够说明它们何处相同、何处不同，以及用什么构造保存这个不同。这种准确辨认与比较的能力，比记住更多性质的名称更值得带到下一次计算和证明中。

\bigskip
\begin{flushright}
R. EverStarine

2026年10月7日
\end{flushright}
